\documentclass[11pt]{article}

\usepackage[margin=1in]{geometry}
\usepackage{amsmath,amssymb,amsthm,mathtools}
\usepackage{booktabs,longtable,array}
\usepackage{algorithm}
\usepackage{algpseudocode}
\usepackage{enumitem}
\usepackage[T1]{fontenc}
\usepackage[utf8]{inputenc}

\newcommand{\Ont}{\mathcal O}
\newcommand{\interp}{\mathcal I}
\newcommand{\DeltaI}{\Delta^{\interp}}
\newcommand{\roleI}[1]{#1^{\interp}}
\newcommand{\conceptI}[1]{#1^{\interp}}
\newcommand{\sub}{\sqsubseteq}
\newcommand{\equivc}{\equiv}
\newcommand{\compose}{\circ}
\newcommand{\occursNeg}{\operatorname{negOccurs}}
\newcommand{\occursPos}{\operatorname{posOccurs}}
\newcommand{\instOf}{\operatorname{instanceOf}}
\newcommand{\init}{\operatorname{init}}
\newcommand{\todo}{\operatorname{todo}}
\newcommand{\subs}{\operatorname{subs}}
\newcommand{\links}{\operatorname{links}}
\newcommand{\hier}{\operatorname{hier}}
\newcommand{\negConjs}{\operatorname{negConjs}}
\newcommand{\negExis}{\operatorname{negExis}}
\newcommand{\conceptIncs}{\operatorname{conceptIncs}}
\newcommand{\roleIncs}{\operatorname{roleIncs}}
\newcommand{\roleComps}{\operatorname{roleComps}}
\newcommand{\superConcepts}{\operatorname{superConcepts}}
\newcommand{\predecessors}{\operatorname{predecessors}}
\newcommand{\successors}{\operatorname{successors}}
\newcommand{\isInitialized}{\operatorname{isInitialized}}
\newcommand{\directSuperConcepts}{\operatorname{directSuperConcepts}}
\newcommand{\equivalentConcepts}{\operatorname{equivalentConcepts}}
\newcommand{\SuperConcept}{\operatorname{SuperConcept}}
\newcommand{\Predecessor}{\operatorname{Predecessor}}
\newcommand{\Successor}{\operatorname{Successor}}
\newcommand{\initializeContext}{\operatorname{initializeContext}}
\newcommand{\process}{\operatorname{process}}
\newcommand{\add}{\operatorname{add}}
\newcommand{\poll}{\operatorname{poll}}

\title{ELK: A Reasoner for OWL EL Ontologies\\Mathematical Content Extracted to LaTeX}
\author{Yevgeny Kazakov, Markus Kr\"otzsch, and Franti\v{s}ek Siman\v{c}\'ik}
\date{}

\begin{document}
\maketitle

\section{Preliminaries}

The vocabulary of $\mathcal{EL}^{+}$ consists of countably infinite sets of atomic roles and atomic concepts. Roles are denoted by $R,S$; concepts by $C,D,E$; atomic concepts by $A,B$. A concept equivalence $C \equivc D$ abbreviates the two inclusions $C \sub D$ and $D \sub C$. An ontology is a finite set of axioms.

Given an ontology $\Ont$, write $\sub^{*}_{\Ont}$ for the smallest reflexive transitive binary relation over roles such that $R \sub^{*}_{\Ont} S$ holds for all $R \sub S \in \Ont$.

A concept $C$ occurs negatively (respectively positively) in an ontology $\Ont$ if $C$ is a syntactic subexpression of $D$ (respectively $E$) for some axiom $D \sub E \in \Ont$.

\subsection{Syntax and Semantics of $\mathcal{EL}^{+}$}

\begin{longtable}{>{\raggedright\arraybackslash}p{0.26\textwidth} >{\raggedright\arraybackslash}p{0.26\textwidth} >{\raggedright\arraybackslash}p{0.38\textwidth}}
\toprule
\textbf{Kind} & \textbf{Syntax} & \textbf{Semantics} \\
\midrule
\endhead
Atomic role & $R$ & $\roleI{R}$ \\
\midrule
Atomic concept & $A$ & $\conceptI{A}$ \\
Top & $\top$ & $\DeltaI$ \\
Conjunction & $C \sqcap D$ & $\conceptI{C} \cap \conceptI{D}$ \\
Existential restriction & $\exists R.C$ & $\{x \mid \exists y : \langle x,y\rangle \in \roleI{R} \land y \in \conceptI{C}\}$ \\
\midrule
Concept inclusion & $C \sub D$ & $\conceptI{C} \subseteq \conceptI{D}$ \\
Role inclusion & $R \sub S$ & $\roleI{R} \subseteq \roleI{S}$ \\
Role composition & $R_1 \compose R_2 \sub S$ & $\langle x,y\rangle \in \roleI{R_1} \land \langle y,z\rangle \in \roleI{R_2} \rightarrow \langle x,z\rangle \in \roleI{S}$ \\
\bottomrule
\end{longtable}

An interpretation $\interp$ consists of a nonempty domain $\DeltaI$ and an interpretation function $\cdot^{\interp}$ assigning to each role $R$ a binary relation $\roleI{R} \subseteq \DeltaI \times \DeltaI$ and to each atomic concept $A$ a set $\conceptI{A} \subseteq \DeltaI$.

$\interp \models \alpha$ means that $\interp$ satisfies axiom $\alpha$. $\interp \models \Ont$ means that $\interp$ satisfies all axioms in $\Ont$. $\Ont \models \alpha$ means that every model of $\Ont$ satisfies $\alpha$. A concept $C$ is subsumed by $D$ with respect to $\Ont$ iff $\Ont \models C \sub D$.

For atomic concepts $A$ and $B$, the subsumption $\Ont \models A \sub B$ is direct if there is no atomic concept $C$ nonequivalent to $A$ and $B$ with respect to $\Ont$ such that $\Ont \models A \sub C$ and $\Ont \models C \sub B$. Classification computes the class taxonomy representing direct subsumptions between equivalence classes of atomic concepts occurring in $\Ont$.

\section{Inference Rules for $\mathcal{EL}^{+}$}

The rules operate with expressions of the forms $C \sub D$, $C \xrightarrow{R} D$, and $\init(C)$. The axiom $C \xrightarrow{R} D$ has the same semantics as $C \sub \exists R.D$ but is used differently by the inference rules. The expression $\init(C)$ initializes derivation of superconcepts for $C$.

\begin{align*}
\textsc{R0} \quad& \frac{\init(C)}{C \sub C}
\end{align*}

\begin{align*}
\textsc{R}_{\top}^{+} \quad&
\frac{\init(C)}{C \sub \top}
\quad \text{if } \top \text{ occurs negatively in } \Ont.
\end{align*}

\begin{align*}
\textsc{R}_{\sqcap}^{-} \quad&
\frac{C \sub D_1 \sqcap D_2}{C \sub D_1 \quad C \sub D_2}
\end{align*}

\begin{align*}
\textsc{R}_{\sqcap}^{+} \quad&
\frac{C \sub D_1 \quad C \sub D_2}{C \sub D_1 \sqcap D_2}
\quad \text{if } D_1 \sqcap D_2 \text{ occurs negatively in } \Ont.
\end{align*}

\begin{align*}
\textsc{R}_{\exists}^{-} \quad&
\frac{C \sub \exists R.D}{\init(D) \quad C \xrightarrow{R} D}
\end{align*}

\begin{align*}
\textsc{R}_{\exists}^{+} \quad&
\frac{E \xrightarrow{R} C \quad C \sub D}{E \sub \exists S.D}
\quad \text{if } \exists S.D \text{ occurs negatively in } \Ont \text{ and } R \sub^{*}_{\Ont} S.
\end{align*}

\begin{align*}
\textsc{R}_{\sub} \quad&
\frac{C \sub D}{C \sub E}
\quad \text{if } D \sub E \in \Ont.
\end{align*}

\begin{align*}
\textsc{R}_{\compose} \quad&
\frac{E \xrightarrow{R_1} C \quad C \xrightarrow{R_2} D}{E \xrightarrow{S} D}
\quad \text{if } S_1 \compose S_2 \sub S \in \Ont,\ R_1 \sub^{*}_{\Ont} S_1,\ R_2 \sub^{*}_{\Ont} S_2.
\end{align*}

The rules are sound. They are complete for classification in the sense that, for each concept $C$ and atomic concept $A$ occurring in $\Ont$, if $\Ont \models C \sub A$, then $C \sub A$ is derivable from $\init(C)$.

\subsection{Example Ontology}

\begin{align}
A &\sub \exists R.(C \sqcap D) \tag{1}\\
B &\equivc A \sqcap \exists S.D \tag{2}\\
\exists S.D &\sub C \tag{3}\\
R &\sub S \tag{4}
\end{align}

To compute all atomic superconcepts of $A$, start with $\init(A)$:

\begin{alignat}{3}
&\init(A) &&\quad \text{initial assumption} &&\tag{5}\\
&A \sub A &&\quad \text{by } \textsc{R0} \text{ to (5)} &&\tag{6}\\
&A \sub \exists R.(C \sqcap D) &&\quad \text{by } \textsc{R}_{\sub} \text{ to (6), using (1)} &&\tag{7}\\
&\init(C \sqcap D) &&\quad \text{by } \textsc{R}_{\exists}^{-} \text{ to (7)} &&\tag{8}\\
&A \xrightarrow{R} C \sqcap D &&\quad \text{by } \textsc{R}_{\exists}^{-} \text{ to (7)} &&\tag{9}\\
&C \sqcap D \sub C \sqcap D &&\quad \text{by } \textsc{R0} \text{ to (8)} &&\tag{10}\\
&C \sqcap D \sub C &&\quad \text{by } \textsc{R}_{\sqcap}^{-} \text{ to (10)} &&\tag{11}\\
&C \sqcap D \sub D &&\quad \text{by } \textsc{R}_{\sqcap}^{-} \text{ to (10)} &&\tag{12}\\
&A \sub \exists S.D &&\quad \text{by } \textsc{R}_{\exists}^{+} \text{ to (9), (12), using (4)} &&\tag{13}\\
&A \sub A \sqcap \exists S.D &&\quad \text{by } \textsc{R}_{\sqcap}^{+} \text{ to (6), (13)} &&\tag{14}\\
&\init(D) &&\quad \text{by } \textsc{R}_{\exists}^{-} \text{ to (13)} &&\tag{15}\\
&A \xrightarrow{S} D &&\quad \text{by } \textsc{R}_{\exists}^{-} \text{ to (13)} &&\tag{16}\\
&A \sub C &&\quad \text{by } \textsc{R}_{\sub} \text{ to (13), using (3)} &&\tag{17}\\
&A \sub B &&\quad \text{by } \textsc{R}_{\sub} \text{ to (14), using (2)} &&\tag{18}\\
&D \sub D &&\quad \text{by } \textsc{R0} \text{ to (15)} &&\tag{19}
\end{alignat}

The derived atomic superconcepts of $A$ include $B$ and $C$, but not $D$.

\section{Indexing}

ELK converts OWL structures into indexed objects for description-logic concepts and roles.

\subsection{Indexed Object Classes}

\begin{center}
\begin{tabular}{lll}
\toprule
\textbf{Class} & \textbf{Fields} & \textbf{Meaning} \\
\midrule
\texttt{IndexedRole} & \texttt{occurs : integer}; \texttt{uri : URI} & Indexed role \\
\texttt{IndexedConcept} & \texttt{negOccurs : integer}; \texttt{posOccurs : integer} & Indexed concept superclass \\
\texttt{IndexedAtomic} & \texttt{uri : URI} & Atomic concept \\
\texttt{IndexedConjunction} & \texttt{firstConj : IndexedConcept}; \texttt{secondConj : IndexedConcept} & Binary conjunction \\
\texttt{IndexedExistential} & \texttt{role : IndexedRole}; \texttt{filler : IndexedConcept} & Existential restriction \\
\bottomrule
\end{tabular}
\end{center}

OWL conjunctions of arbitrary arity are converted into nested binary conjunctions; for example, $A \sqcap B \sqcap C \sqcap D$ is represented as $((A \sqcap B) \sqcap C) \sqcap D$.

\subsection{Index Tables}

\begin{align*}
\negConjs &= \{\langle C,D,C \sqcap D\rangle \mid C \sqcap D \text{ occurs negatively in } \Ont\},\\
\negExis &= \{\langle R,C,\exists R.C\rangle \mid \exists R.C \text{ occurs negatively in } \Ont\},\\
\conceptIncs &= \{\langle C,D\rangle \mid C \sub D \in \Ont\},\\
\roleIncs &= \{\langle R,S\rangle \mid R \sub S \in \Ont\},\\
\roleComps &= \{\langle R_1,R_2,S\rangle \mid R_1 \compose R_2 \sub S \in \Ont\}.
\end{align*}

\subsection{Index for the Example Ontology}

\begin{center}
\begin{tabular}{lrrrrrrrrr}
\toprule
& $\top$ & $A$ & $B$ & $C$ & $D$ & $C\sqcap D$ & $\exists R.(C\sqcap D)$ & $\exists S.D$ & $A\sqcap \exists S.D$ \\
\midrule
$\occursNeg$ & 0 & 2 & 1 & 0 & 2 & 0 & 0 & 2 & 1 \\
$\occursPos$ & 0 & 1 & 1 & 2 & 2 & 1 & 1 & 1 & 1 \\
\bottomrule
\end{tabular}
\end{center}

\begin{center}
\begin{tabular}{lcc}
\toprule
\textbf{IndexedRole} & $R$ & $S$ \\
\midrule
\texttt{occurs} & 2 & 4 \\
\bottomrule
\end{tabular}
\end{center}

For the example ontology:
\begin{align*}
\negConjs &= \{\langle A,\exists S.D,A\sqcap \exists S.D\rangle\},\\
\negExis &= \{\langle S,D,\exists S.D\rangle\},\\
\conceptIncs &= \{\langle A,\exists R.(C\sqcap D)\rangle,
\langle B,A\sqcap \exists S.D\rangle,
\langle A\sqcap \exists S.D,B\rangle,
\langle \exists S.D,C\rangle\},\\
\roleIncs &= \{\langle R,S\rangle\},\\
\roleComps &= \emptyset.
\end{align*}

For deletion of $B \equivc A \sqcap \exists S.D$, tuples $\langle B,A\sqcap\exists S.D\rangle$ and $\langle A\sqcap\exists S.D,B\rangle$ are removed from $\conceptIncs$. Negative and positive counts of $B$, $A\sqcap\exists S.D$, $A$, and $\exists S.D$ are decremented; the occurrence count of $S$ is decremented by $2$. Since $\occursNeg(A\sqcap\exists S.D)$ drops to zero, $\langle A,\exists S.D,A\sqcap\exists S.D\rangle$ is removed from $\negConjs$. Since $\occursNeg(\exists S.D)$ remains nonzero, $\negExis$ is unchanged.

\section{Saturation of Roles}

Precompute the reflexive transitive closure $\sub^{*}_{\Ont}$ of role inclusion axioms and store it in:
\begin{align*}
\hier = \{\langle R,S\rangle \mid R \sub^{*}_{\Ont} S\}.
\end{align*}

For the example ontology:
\begin{align*}
\hier = \{\langle R,R\rangle,\langle R,S\rangle,\langle S,S\rangle\}.
\end{align*}

\section{Saturation of Concepts}

The three types of derivable axioms are represented internally as:

\begin{center}
\begin{tabular}{ll}
\toprule
\textbf{Axiom} & \textbf{Internal representation} \\
\midrule
$\init(C)$ & $C \in \init$ \\
$C \sub D$ & $\langle C,D\rangle \in \subs$ \\
$C \xrightarrow{R} D$ & $\langle C,R,D\rangle \in \links$ \\
\bottomrule
\end{tabular}
\end{center}

\subsection{Compiled Inference Rules}

Here $T(x)$ abbreviates $x \in T$.

\begin{align*}
\textsc{R0}:&\quad \text{If } \init(C), \text{ then } \subs(C,C).\\
\textsc{R}_{\top}^{+}:&\quad \text{If } \init(C) \land \top.\occursNeg > 0, \text{ then } \subs(C,\top).\\
\textsc{R}_{\sqcap}^{-}:&\quad \text{If } \subs(C,D) \land D \instOf \texttt{IndexedConjunction},\\
&\quad \text{then } \subs(C,D.\texttt{firstConj}) \text{ and } \subs(C,D.\texttt{secondConj}).\\
\textsc{R}_{\sqcap}^{+}:&\quad \text{If } \subs(C,D_1) \land \subs(C,D_2) \land \negConjs(D_1,D_2,D),\\
&\quad \text{then } \subs(C,D).\\
\textsc{R}_{\exists}^{-}:&\quad \text{If } \subs(C,D) \land D \instOf \texttt{IndexedExistential},\\
&\quad \text{then } \init(D.\texttt{filler}) \text{ and } \links(C,D.\texttt{role},D.\texttt{filler}).\\
\textsc{R}_{\exists}^{+}:&\quad \text{If } \links(E,R,C) \land \subs(C,D) \land \negExis(S,D,F) \land \hier(R,S),\\
&\quad \text{then } \subs(E,F).\\
\textsc{R}_{\sub}:&\quad \text{If } \subs(C,D) \land \conceptIncs(D,E), \text{ then } \subs(C,E).\\
\textsc{R}_{\compose}:&\quad \text{If } \links(E,R_1,C) \land \links(C,R_2,D) \land \roleComps(S_1,S_2,S)\\
&\quad \land\ \hier(R_1,S_1) \land \hier(R_2,S_2), \text{ then } \links(E,S,D).
\end{align*}

\begin{algorithm}
\caption{Main body of the saturation algorithm}
\begin{algorithmic}[1]
\ForAll{\texttt{IndexedAtomic} $A$}
    \State $\todo.\add(\init(A))$
\EndFor
\While{$\todo \neq \emptyset$}
    \State $\mathit{axiom} \gets \todo.\poll()$
    \State $\process(\mathit{axiom})$
\EndWhile
\end{algorithmic}
\end{algorithm}

\begin{algorithm}
\caption{$\process(\mathit{axiom})$}
\begin{algorithmic}[1]
\Procedure{process}{$\init(C)$}
    \If{$\init.\add(C)$}
        \State $\todo.\add(C \sub C)$ \Comment{rule $\textsc{R0}$}
        \If{$\top.\occursNeg > 0$}
            \State $\todo.\add(C \sub \top)$ \Comment{rule $\textsc{R}_{\top}^{+}$}
        \EndIf
    \EndIf
\EndProcedure
\Procedure{process}{$C \sub D$}
    \If{$\subs.\add(\langle C,D\rangle)$}
        \If{$D \instOf \texttt{IndexedConjunction}$}
            \State $\todo.\add(C \sub D.\texttt{firstConj})$
            \State $\todo.\add(C \sub D.\texttt{secondConj})$ \Comment{rule $\textsc{R}_{\sqcap}^{-}$}
        \EndIf
        \If{$D \instOf \texttt{IndexedExistential}$}
            \State $\todo.\add(\init(D.\texttt{filler}))$
            \State $\todo.\add(C \xrightarrow{D.\texttt{role}} D.\texttt{filler})$ \Comment{rule $\textsc{R}_{\exists}^{-}$}
        \EndIf
        \ForAll{$D_2,E$ with $\subs(C,D_2) \land \negConjs(D,D_2,E)$}
            \State $\todo.\add(C \sub E)$ \Comment{rule $\textsc{R}_{\sqcap}^{+}$}
        \EndFor
        \ForAll{$D_1,E$ with $\subs(C,D_1) \land \negConjs(D_1,D,E)$}
            \State $\todo.\add(C \sub E)$ \Comment{rule $\textsc{R}_{\sqcap}^{+}$}
        \EndFor
        \ForAll{$E,F,R,S$ with $\links(E,R,C) \land \negExis(S,D,F) \land \hier(R,S)$}
            \State $\todo.\add(E \sub F)$ \Comment{rule $\textsc{R}_{\exists}^{+}$}
        \EndFor
        \ForAll{$E$ with $\conceptIncs(D,E)$}
            \State $\todo.\add(C \sub E)$ \Comment{rule $\textsc{R}_{\sub}$}
        \EndFor
    \EndIf
\EndProcedure
\Procedure{process}{$E \xrightarrow{R} C$}
    \If{$\links.\add(\langle E,R,C\rangle)$}
        \ForAll{$D,F,S$ with $\subs(C,D) \land \negExis(S,D,F) \land \hier(R,S)$}
            \State $\todo.\add(E \sub F)$ \Comment{rule $\textsc{R}_{\exists}^{+}$}
        \EndFor
        \ForAll{$D,R_2,S_1,S_2,S$ with $\links(C,R_2,D) \land \roleComps(S_1,S_2,S) \land \hier(R,S_1) \land \hier(R_2,S_2)$}
            \State $\todo.\add(E \xrightarrow{S} D)$ \Comment{rule $\textsc{R}_{\compose}$}
        \EndFor
        \ForAll{$D,R_1,S_1,S_2,S$ with $\links(D,R_1,E) \land \roleComps(S_1,S_2,S) \land \hier(R_1,S_1) \land \hier(R,S_2)$}
            \State $\todo.\add(D \xrightarrow{S} C)$ \Comment{rule $\textsc{R}_{\compose}$}
        \EndFor
    \EndIf
\EndProcedure
\end{algorithmic}
\end{algorithm}

\subsection{Saturation for the Example Ontology}

\begin{align*}
\init &= \{A,B,C,D,C\sqcap D\},\\
\subs &= \{\langle A,A\rangle,\langle A,\exists R.(C\sqcap D)\rangle,\langle A,\exists S.D\rangle,
\langle A,A\sqcap\exists S.D\rangle,\langle A,C\rangle,\langle A,B\rangle,\\
&\qquad \langle B,B\rangle,\langle B,A\sqcap\exists S.D\rangle,\langle B,A\rangle,
\langle B,\exists S.D\rangle,\langle B,\exists R.(C\sqcap D)\rangle,\langle B,C\rangle,\\
&\qquad \langle C,C\rangle,\langle D,D\rangle,\langle C\sqcap D,C\sqcap D\rangle,
\langle C\sqcap D,C\rangle,\langle C\sqcap D,D\rangle\},\\
\links &= \{\langle A,R,C\sqcap D\rangle,\langle A,S,D\rangle,\langle B,S,D\rangle,\langle B,R,C\sqcap D\rangle\}.
\end{align*}

\section{Taxonomy Construction}

After concept saturation, discard all subsumptions involving non-atomic concepts. The remaining atomic subsumptions are transitively reduced.

\begin{algorithm}
\caption{Naive transitive reduction}
\begin{algorithmic}[1]
\ForAll{$C$ with $\subs(A,C)$}
    \State $\mathit{isDirect} \gets \textbf{true}$
    \ForAll{$B$ with $\subs(A,B)$}
        \If{$B \neq C$ and $\subs(B,C)$}
            \State $\mathit{isDirect} \gets \textbf{false}$
        \EndIf
    \EndFor
    \If{$\mathit{isDirect}$}
        \State $A.\directSuperConcepts.\add(C)$
    \EndIf
\EndFor
\end{algorithmic}
\end{algorithm}

\begin{algorithm}
\caption{Better transitive reduction}
\begin{algorithmic}[1]
\ForAll{$C$ with $\subs(A,C)$}
    \If{$\subs(C,A)$}
        \State $A.\equivalentConcepts.\add(C)$
    \Else
        \State $\mathit{isDirect} \gets \textbf{true}$
        \ForAll{$B \in A.\directSuperConcepts$}
            \If{$\subs(B,C)$}
                \State $\mathit{isDirect} \gets \textbf{false}$
                \State \textbf{break}
            \EndIf
            \If{$\subs(C,B)$}
                \State $A.\directSuperConcepts.\operatorname{remove}(B)$
            \EndIf
        \EndFor
        \If{$\mathit{isDirect}$}
            \State $A.\directSuperConcepts.\add(C)$
        \EndIf
    \EndIf
\EndFor
\end{algorithmic}
\end{algorithm}

For the example ontology, projecting to atomic concepts leaves:
\begin{align*}
\mathrm{Subs} = \{&\langle A,A\rangle,\langle A,B\rangle,\langle A,C\rangle,
\langle B,A\rangle,\langle B,B\rangle,\langle B,C\rangle,\langle C,C\rangle,\langle D,D\rangle\}.
\end{align*}

Algorithm 4 computes:
\begin{align*}
A.\equivalentConcepts &= \{A,B\}, & A.\directSuperConcepts &= \{C\},\\
B.\equivalentConcepts &= \{A,B\}, & B.\directSuperConcepts &= \{C\},\\
C.\equivalentConcepts &= \{C\}, & C.\directSuperConcepts &= \emptyset,\\
D.\equivalentConcepts &= \{D\}, & D.\directSuperConcepts &= \emptyset.
\end{align*}

The taxonomy has top node $\top$, then $C$ and $D$ beneath $\top$, then the equivalence class $\{A,B\}$ beneath $C$, and bottom node $\bot$ beneath all leaf classes.

\section{Join Iteration Optimizations}

For $\textsc{R}_{\sqcap}^{+}$, when processing $C \sub D$ as first premise, line 14 of Algorithm 2 iterates over all $D_2,E$ such that:
\begin{align*}
\subs(C,D_2) \land \negConjs(D,D_2,E).
\end{align*}

A cached implication optimization stores partially matched conjunctions. When processing $C \sub B_i$, if $C \sub A$ has not been derived, store:
\begin{align*}
C \sub (A \to A \sqcap B_i)
\end{align*}
represented as a triple:
\begin{align*}
\langle C,A,A\sqcap B_i\rangle \in \operatorname{impl}.
\end{align*}
When processing $C \sub A$, retrieve all $E$ with $\operatorname{impl}(C,A,E)$ and derive $C \sub E$.

For $\textsc{R}_{\exists}^{+}$, ignoring role hierarchy, applying the rule to $E \xrightarrow{R} C$ requires all $D,F$ such that:
\begin{align*}
\subs(C,D) \land \negExis(R,D,F).
\end{align*}

A propagation optimization precomputes the join. When processing $C \sub D$, derive propagation axioms:
\begin{align*}
C \sub (R \to F)
\end{align*}
for each $R,F$ such that $\negExis(R,D,F)$, represented as:
\begin{align*}
\langle C,R,F\rangle \in \operatorname{prop}.
\end{align*}
When processing $E \xrightarrow{R} C$, retrieve all $F$ with $\operatorname{prop}(C,R,F)$ and derive $E \sub F$.

This corresponds to the rule:
\begin{align*}
\textsc{RH} \quad \frac{E \xrightarrow{R} C \quad C \sub (R \to F)}{E \sub F}.
\end{align*}

With role hierarchy, one possible form is:
\begin{align*}
\textsc{RH} \quad \frac{E \xrightarrow{R} C \quad C \sub (S \to F)}{E \sub F}
\quad \text{if } R \sub^{*}_{\Ont} S.
\end{align*}

Alternative expansions under the role hierarchy:
\begin{align*}
\frac{E \xrightarrow{R} C}{E \xrightarrow{S} C}
\quad \text{if } R \sub^{*}_{\Ont} S,
\end{align*}
\begin{align*}
\frac{C \sub (S \to F)}{C \sub (R \to F)}
\quad \text{if } R \sub^{*}_{\Ont} S.
\end{align*}

\section{Concurrent Saturation}

Axioms are assigned to contexts. There is one context for each initialized concept $C$. Axioms of the form $\init(C)$, $C \sub D$, or $E \xrightarrow{R} C$ are assigned to the context of $C$. In the presence of role composition axioms, $E \xrightarrow{R} C$ is also assigned to the context of $E$. Whenever a binary inference rule $\textsc{R}_{\sqcap}^{+}$, $\textsc{R}_{\exists}^{+}$, or $\textsc{R}_{\compose}$ is applicable to a pair of premises, the two premises belong to the same context.

Internal representation in the refined saturation algorithm:

\begin{center}
\begin{tabular}{lll}
\toprule
\textbf{Axiom} & \textbf{To-do representation} & \textbf{Processed representation} \\
\midrule
$\init(C)$ & -- & $C.\isInitialized = \textbf{true}$ \\
$C \sub D$ & $\SuperConcept(D) \in C.\todo$ & $D \in C.\superConcepts$ \\
$E \xrightarrow{R} C$ & $\Predecessor(R,E) \in C.\todo$ & $\langle R,E\rangle \in C.\predecessors$ \\
$E \xrightarrow{R} C$ & $\Successor(R,C) \in E.\todo$ & $\langle R,C\rangle \in E.\successors$ \\
\bottomrule
\end{tabular}
\end{center}

Storing $E \xrightarrow{R} C$ in $E.\successors$ is needed only in the presence of role compositions for rule $\textsc{R}_{\compose}$.

\begin{algorithm}
\caption{Context-local processing $C.\process(\mathit{axiom})$}
\begin{algorithmic}[1]
\Procedure{initializeContext}{$C$}
    \If{$C.\isInitialized = \textbf{false}$}
        \State $C.\isInitialized \gets \textbf{true}$
        \State $C.\todo.\add(\SuperConcept(C))$ \Comment{rule $\textsc{R0}$}
        \If{$\top.\occursNeg > 0$}
            \State $C.\todo.\add(\SuperConcept(\top))$ \Comment{rule $\textsc{R}_{\top}^{+}$}
        \EndIf
    \EndIf
\EndProcedure
\Procedure{process}{$\SuperConcept(D)$}
    \If{$C.\superConcepts.\add(D)$}
        \If{$D \instOf \texttt{IndexedConjunction}$}
            \State $C.\todo.\add(\SuperConcept(D.\texttt{firstConj}))$
            \State $C.\todo.\add(\SuperConcept(D.\texttt{secondConj}))$ \Comment{rule $\textsc{R}_{\sqcap}^{-}$}
        \EndIf
        \If{$D \instOf \texttt{IndexedExistential}$}
            \State $\initializeContext(D.\texttt{filler})$
            \State $D.\texttt{filler}.\todo.\add(\Predecessor(D.\texttt{role},C))$ \Comment{rule $\textsc{R}_{\exists}^{-}$}
        \EndIf
        \ForAll{$D_2,E$ with $C.\superConcepts(D_2) \land \negConjs(D,D_2,E)$}
            \State $C.\todo.\add(\SuperConcept(E))$ \Comment{rule $\textsc{R}_{\sqcap}^{+}$}
        \EndFor
        \ForAll{$D_1,E$ with $C.\superConcepts(D_1) \land \negConjs(D_1,D,E)$}
            \State $C.\todo.\add(\SuperConcept(E))$ \Comment{rule $\textsc{R}_{\sqcap}^{+}$}
        \EndFor
        \ForAll{$E,F,R,S$ with $C.\predecessors(R,E) \land \negExis(S,D,F) \land \hier(R,S)$}
            \State $E.\todo.\add(\SuperConcept(F))$ \Comment{rule $\textsc{R}_{\exists}^{+}$}
        \EndFor
        \ForAll{$E$ with $\conceptIncs(D,E)$}
            \State $C.\todo.\add(\SuperConcept(E))$ \Comment{rule $\textsc{R}_{\sub}$}
        \EndFor
    \EndIf
\EndProcedure
\Procedure{process}{$\Predecessor(R,E)$}
    \If{$C.\predecessors.\add(\langle R,E\rangle)$}
        \ForAll{$D,F,S$ with $C.\superConcepts(D) \land \negExis(S,D,F) \land \hier(R,S)$}
            \State $E.\todo.\add(\SuperConcept(F))$ \Comment{rule $\textsc{R}_{\exists}^{+}$}
        \EndFor
        \ForAll{$D,R_2,S_1,S_2,S$ with $C.\successors(R_2,D) \land \roleComps(S_1,S_2,S) \land \hier(R,S_1) \land \hier(R_2,S_2)$}
            \State $D.\todo.\add(\Predecessor(S,E))$ \Comment{rule $\textsc{R}_{\compose}$}
        \EndFor
        \State $E.\todo.\add(\Successor(R,C))$ \Comment{only needed with role composition axioms}
    \EndIf
\EndProcedure
\Procedure{process}{$\Successor(R,D)$}
    \If{$C.\successors.\add(\langle R,D\rangle)$}
        \ForAll{$E,R_1,S_1,S_2,S$ with $C.\predecessors(R_1,E) \land \roleComps(S_1,S_2,S) \land \hier(R_1,S_1) \land \hier(R,S_2)$}
            \State $D.\todo.\add(\Predecessor(S,E))$ \Comment{rule $\textsc{R}_{\compose}$}
        \EndFor
    \EndIf
\EndProcedure
\end{algorithmic}
\end{algorithm}

\section{Composition-Decomposition Optimizations}

The system can omit certain applications of decomposition rules without losing completeness.

\begin{description}[leftmargin=2.5em]
\item[$O_{\sqcap}$] It is not necessary to apply $\textsc{R}_{\sqcap}^{-}$ to a subsumption $C \sub D_1 \sqcap D_2$ that was derived by $\textsc{R}_{\sqcap}^{+}$.
\item[$O_{\exists}$] It is not necessary to apply $\textsc{R}_{\exists}^{-}$ to a subsumption $C \sub \exists S.E$ that was derived by $\textsc{R}_{\exists}^{+}$.
\end{description}

For $O_{\sqcap}$, if $C \sub D_1 \sqcap D_2$ was derived by $\textsc{R}_{\sqcap}^{+}$, then both $C \sub D_1$ and $C \sub D_2$ were derived earlier.

For $O_{\exists}$, if $C \sub \exists S.E$ was derived by $\textsc{R}_{\exists}^{+}$, then some $C \xrightarrow{R} D$ with $R \sub^{*}_{\Ont} S$ and $D \sub E$ was derived earlier. The conclusion $C \xrightarrow{S} E$ of $\textsc{R}_{\exists}^{-}$ is therefore redundant for satisfying the existential restriction.

In the example derivation, applying $\textsc{R}_{\exists}^{-}$ to axiom (13) is redundant under $O_{\exists}$; this avoids deriving axioms (15), (16), and (19), and avoids introducing the context of $D$.

\end{document}
